% !TeX root = ../print.tex
\section{EasyUC Secure Message Communication}

\newcommand{\addr}{\mathsf{addr}}
\newcommand{\port}{\mathsf{port}}
\newcommand{\dir}{\mathsf{dir}}
\newcommand{\advmsg}{\mathsf{adv}}
\renewcommand{\Exec}{\mathsf{Exec}}
\newcommand{\MI}{\mathsf{MI}}
\renewcommand{\Forw}{\mathsf{Forw}}
\newcommand{\KEReal}{\mathsf{KEReal}}
\newcommand{\KEIdeal}{\mathsf{KEIdeal}}
\newcommand{\KESim}{\mathsf{KESim}}
\newcommand{\SMCReal}{\mathsf{SMCReal}}
\newcommand{\SMCIdeal}{\mathsf{SMCIdeal}}
\newcommand{\SMCSim}{\mathsf{SMCSim}}
\newcommand{\SMCSimComp}{\mathsf{SMCSimComp}}
\newcommand{\CompEnv}{\mathsf{CompEnv}}
\newcommand{\DDHOne}{\mathsf{DDH}_1}
\newcommand{\DDHTwo}{\mathsf{DDH}_2}
\newcommand{\Adv}{\mathsf{Adv}}
\newcommand{\Env}{\mathsf{Env}}
\newcommand{\Sim}{\mathsf{Sim}}
\newcommand{\Key}{\mathsf{key}}
\newcommand{\Exp}{\mathsf{exp}}
\newcommand{\Text}{\mathsf{text}}
\newcommand{\Univ}{\mathsf{univ}}
\newcommand{\dexp}{\mathsf{dexp}}
\newcommand{\inj}{\mathsf{inj}}
\renewcommand{\proj}{\mathsf{proj}}
\newcommand{\oget}{\mathsf{oget}}
\newcommand{\kinv}{\mathsf{kinv}}
\newcommand{\gid}{\mathsf{kid}}
\newcommand{\mul}{\mathbin{\hat{\hat{}}}}
\newcommand{\pow}{\mathbin{\hat{}}}
\newcommand{\pt}{\mathsf{pt}}
\newcommand{\true}{\mathsf{true}}
\newcommand{\false}{\mathsf{false}}
\newcommand{\res}{\mathsf{res}}

This section states a self-contained proof of UC security for the EasyUC secure
message communication case study.  The model is the restricted EasyUC model:
all machines are statically created modules, addresses are hierarchical, and an
experiment consists of an interface containing a functionality and an adversary,
interacting with an environment.  A real functionality may contain
sub-functionalities; the address of the \(i\)-th sub-functionality of address
\(\alpha\) is written \(\alpha i\).

\subsection{Preliminaries}

\begin{definition}[Ports and messages]
A port is a pair \((\alpha,i)\), where \(\alpha\in\addr\) is a machine address
and \(i\in\mathbb{N}\) is a port index.  Messages are of two kinds.
A direct message is written
\[
  p \xrightarrow{x} q ,
\]
where \(p,q\) are ports and \(x\in\Univ\).  An adversarial message is written
\[
  p \xRightarrow{x} q .
\]
The interface routes direct messages to the functionality side and adversarial
messages to the adversary side, subject to the input guard.  Messages whose
kind, source, destination, or payload shape is not accepted by the receiving
machine are rejected and produce no output.
\end{definition}

\begin{definition}[Execution experiment]
For a functionality \(F\), an adversary \(A\), an environment \(Z\), addresses
\(\alpha,\beta\), and an input guard \(G\), let
\[
  \Exec(F,A,Z;\alpha,\beta,G,n)
\]
be the Bernoulli random variable obtained by running the EasyUC controller at
security parameter \(n\), with interface \(\MI(F,A)\), functionality address
\(\alpha\), adversary address \(\beta\), input guard \(G\), and environment
\(Z\), until \(Z\) returns its final bit.  The probability that the output bit
is \(1\) is written
\[
  \Pr[\Exec(F,A,Z;\alpha,\beta,G,n)=1].
\]
All probabilities are over the coins of all randomized machines.
\end{definition}

\begin{definition}[Computational emulation with concrete bound]
For functionalities \(F_0,F_1\), adversary transformer \(S\), and bound
\(\varepsilon\), write
\[
  F_0 \preceq_{\varepsilon,S} F_1
\]
if for every PPT adversary \(A\), every PPT environment \(Z\), every valid
address pair \(\alpha,\beta\), every allowed input guard \(G\), and every
security parameter \(n\),
\[
\begin{aligned}
&\left|
  \Pr[\Exec(F_0,A,Z;\alpha,\beta,G,n)=1]
 -\Pr[\Exec(F_1,S(A),Z;\alpha,\beta,G,n)=1]
 \right|  \\
&\hspace{4em}\leq \varepsilon(A,Z,\alpha,\beta,G,n).
\end{aligned}
\]
If \(\varepsilon\) is negligible for all PPT \(A,Z\), this is computational
emulation.
\end{definition}

\begin{definition}[Algebraic setting]
Let \(\Exp\) be a finite exponent type with uniform distribution \(\dexp\).
Let \((\Key,\mul,\gid,\kinv)\) be a finite group.  Let \(g\in\Key\), and let
\(\pow:\Key\times\Exp\to\Key\) be exponentiation.  We assume:
\[
  (g\pow q_1)\pow q_2 = g\pow(q_1q_2)
  \quad\text{and}\quad
  q_1q_2=q_2q_1.
\]
There is an injection \(\inj:\Text\to\Key\) and a partial projection
\(\proj:\Key\to\mathsf{option}(\Text)\) satisfying
\[
  \proj(\inj(t))=\mathsf{Some}(t)
  \quad\text{for all }t\in\Text.
\]
The one-time-pad encryption of \(t\) under key \(k\) is
\[
  \mathsf{otp}(k,t)=\inj(t)\mul k,
\]
and decryption is
\[
  \mathsf{otpdec}(k,e)=\oget(\proj(e\mul\kinv(k))).
\]
\end{definition}

\begin{lemma}[One-time-pad correctness]
For all \(t\in\Text\) and \(k\in\Key\),
\[
  \mathsf{otpdec}(k,\mathsf{otp}(k,t))=t.
\]
\end{lemma}

\begin{proof}
By definition,
\[
\begin{aligned}
 \mathsf{otpdec}(k,\mathsf{otp}(k,t))
 &= \oget(\proj((\inj(t)\mul k)\mul\kinv(k)))       \\
 &= \oget(\proj(\inj(t)\mul(k\mul\kinv(k))))        \\
 &= \oget(\proj(\inj(t)\mul\gid))                   \\
 &= \oget(\proj(\inj(t)))                           \\
 &= \oget(\mathsf{Some}(t))                         \\
 &= t.
\end{aligned}
\]
Only associativity, the inverse law, the identity law, and
\(\proj(\inj(t))=\mathsf{Some}(t)\) are used.
\end{proof}

\begin{lemma}[One-time-pad masking]
Let \(K\) be uniformly distributed over \(\Key\).  For every fixed
\(t\in\Text\), the random variable \(\inj(t)\mul K\) is uniformly distributed
over \(\Key\).  Consequently, for all \(t_0,t_1\in\Text\),
\[
  \inj(t_0)\mul K \equiv \inj(t_1)\mul K
\]
as distributions.
\end{lemma}

\begin{proof}
Fix \(t\).  The map
\[
  \varphi_t:\Key\to\Key,\qquad \varphi_t(k)=\inj(t)\mul k
\]
is a bijection.  Its inverse is
\[
  \psi_t(x)=\kinv(\inj(t))\mul x.
\]
Indeed,
\[
  \psi_t(\varphi_t(k))
  =\kinv(\inj(t))\mul(\inj(t)\mul k)
  =(\kinv(\inj(t))\mul\inj(t))\mul k
  =\gid\mul k
  =k,
\]
and similarly \(\varphi_t(\psi_t(x))=x\).  A bijection preserves the uniform
distribution on a finite set.  Thus \(\inj(t)\mul K\) is uniform over
\(\Key\), independently of \(t\).  The final distributional equality follows.
\end{proof}

\subsection{Functionalities and protocols}

\begin{definition}[The forwarding functionality \(\Forw\)]
\(\Forw\) is one-shot.  It has a functionality address \(\alpha\), an adversary
address \(\beta\), one port index \(1\), and three states:
\[
  \mathsf{FState}_1,\qquad
  \mathsf{FState}_2(\pt_1,\pt_2,u),\qquad
  \mathsf{FState}_3.
\]
Its input guard accepts direct and adversarial messages to port \((\alpha,1)\)
only.

\begin{enumerate}
\item In state \(\mathsf{FState}_1\), on a direct message
\[
  \pt_1 \xrightarrow{(\pt_2,u)} (\alpha,1),
\]
where \(\pt_1,\pt_2\) are external to \(\alpha\) and \(\beta\), it outputs
\[
  (\alpha,1) \xRightarrow{(\pt_1,\pt_2,u)} (\beta,1)
\]
and moves to \(\mathsf{FState}_2(\pt_1,\pt_2,u)\).

\item In state \(\mathsf{FState}_2(\pt_1,\pt_2,u)\), on an adversarial approval
message from \((\beta,1)\) to \((\alpha,1)\), it outputs
\[
  (\alpha,1) \xrightarrow{(\pt_1,u)} \pt_2
\]
and moves to \(\mathsf{FState}_3\).

\item In state \(\mathsf{FState}_3\), every message is rejected.
\end{enumerate}
Every message not matching the current state is rejected and produces no
output.
\end{definition}

\begin{lemma}[Trace theorem for \(\Forw\)]
For every initial state of \(\Forw\), every input message \(m\), and every
output/state pair \((o,s')\), the transition relation of \(\Forw\) satisfies
exactly the following clauses:
\[
\begin{array}{ll}
\mathsf{FState}_1
 \xrightarrow{\;\pt_1 \xrightarrow{(\pt_2,u)} (\alpha,1)\;}
 \big((\alpha,1)\xRightarrow{(\pt_1,\pt_2,u)}(\beta,1),
       \mathsf{FState}_2(\pt_1,\pt_2,u)\big),\\[0.4em]
\mathsf{FState}_2(\pt_1,\pt_2,u)
 \xrightarrow{\;(\beta,1)\xRightarrow{\mathsf{ok}}(\alpha,1)\;}
 \big((\alpha,1)\xrightarrow{(\pt_1,u)}\pt_2,\mathsf{FState}_3\big),\\[0.4em]
\mathsf{FState}_3 \xrightarrow{\;m\;}(\bot,\mathsf{FState}_3),
\end{array}
\]
and all non-matching cases return \((\bot,s)\), where \(s\) is the prior state.
\end{lemma}

\begin{proof}
The proof is by case analysis on the current state and then on the syntactic
form of the incoming message.  In \(\mathsf{FState}_1\), the only accepted
message is a direct message to \((\alpha,1)\) with payload \((\pt_2,u)\) and
valid external source \(\pt_1\); unfolding the transition function gives the
first displayed transition and no other output.  In
\(\mathsf{FState}_2(\pt_1,\pt_2,u)\), the only accepted message is an
adversarial approval from \((\beta,1)\) to \((\alpha,1)\); unfolding gives the
second displayed transition.  In \(\mathsf{FState}_3\), the definition rejects
all messages.  The non-matching cases are exactly the remaining constructors
of the message type and the failed guard checks, each of which returns
\(\bot\) and leaves the state unchanged.
\end{proof}

\begin{definition}[The ideal secure-message functionality \(\SMCIdeal\)]
\(\SMCIdeal\) has address \(\alpha\), adversary address \(\beta\), port indices
\(1,2,3\), and three states:
\[
  \mathsf{SState}_1,\qquad
  \mathsf{SState}_2(\pt_1,\pt_2,t),\qquad
  \mathsf{SState}_3.
\]
Port \(1\) is the sender-side direct port, port \(2\) is the receiver-side
direct output port, and port \(3\) communicates with the simulator.  The input
guard accepts direct messages to port \(1\) and adversarial messages to port
\(3\).

\begin{enumerate}
\item In state \(\mathsf{SState}_1\), on
\[
  \pt_1 \xrightarrow{(\pt_2,t)}(\alpha,1),
\]
where \(\pt_1,\pt_2\) are external to \(\alpha,\beta\), it outputs
\[
  (\alpha,3)\xRightarrow{(\pt_1,\pt_2)}(\beta,3)
\]
and moves to \(\mathsf{SState}_2(\pt_1,\pt_2,t)\).  The plaintext \(t\) is not
sent to the adversary.

\item In state \(\mathsf{SState}_2(\pt_1,\pt_2,t)\), on adversarial approval
from \((\beta,3)\) to \((\alpha,3)\), it outputs
\[
  (\alpha,2)\xrightarrow{(\pt_1,t)}\pt_2
\]
and moves to \(\mathsf{SState}_3\).

\item In state \(\mathsf{SState}_3\), all messages are rejected.
\end{enumerate}
\end{definition}

\begin{definition}[The ideal key-exchange functionality \(\KEIdeal\)]
\(\KEIdeal\) has address \(\alpha\), adversary address \(\beta\), port indices
\(1,2,3\), and five states:
\[
\begin{array}{c}
  \mathsf{KState}_1,\quad
  \mathsf{KState}_2(\pt_1,\pt_2),\quad
  \mathsf{KState}_3(\pt_1,\pt_2,k),\\
  \mathsf{KState}_4(\pt_1,\pt_2,k),\quad
  \mathsf{KState}_5.
\end{array}
\]
The input guard accepts direct messages to ports \(1,2\) and adversarial
messages to port \(3\).

\begin{enumerate}
\item In \(\mathsf{KState}_1\), on direct input
\[
  \pt_1\xrightarrow{\pt_2}(\alpha,1),
\]
it sends
\[
  (\alpha,3)\xRightarrow{(\pt_1,\pt_2)}(\beta,2)
\]
and moves to \(\mathsf{KState}_2(\pt_1,\pt_2)\).

\item In \(\mathsf{KState}_2(\pt_1,\pt_2)\), on approval from \((\beta,2)\) to
\((\alpha,3)\), it samples \(q\leftarrow\dexp\), sets \(k=g\pow q\), sends
\[
  (\alpha,2)\xrightarrow{(\pt_1,k)}\pt_2,
\]
and moves to \(\mathsf{KState}_3(\pt_1,\pt_2,k)\).

\item In \(\mathsf{KState}_3(\pt_1,\pt_2,k)\), on direct input from
\(\pt_2\) to \((\alpha,2)\), it sends
\[
  (\alpha,3)\xRightarrow{\star}(\beta,2)
\]
and moves to \(\mathsf{KState}_4(\pt_1,\pt_2,k)\).

\item In \(\mathsf{KState}_4(\pt_1,\pt_2,k)\), on approval from \((\beta,2)\)
to \((\alpha,3)\), it sends
\[
  (\alpha,1)\xrightarrow{k}\pt_1
\]
and moves to \(\mathsf{KState}_5\).

\item In \(\mathsf{KState}_5\), all messages are rejected.
\end{enumerate}
\end{definition}

\begin{definition}[The real key-exchange functionality \(\KEReal\)]
\(\KEReal\) has address \(\alpha\), adversary address \(\beta\), two
sub-functionalities \(\Forw_1\) at \(\alpha1\) and \(\Forw_2\) at \(\alpha2\),
and party ports \(1,2,3,4\).  Ports \(1,3\) belong to Party~1; ports \(2,4\)
belong to Party~2.  Its input guard accepts direct messages to ports \(1,2\)
and adversarial messages addressed to \(\alpha1,\alpha2\).

Party~1 has three states:
\[
  \mathsf{P1}_1,\qquad
  \mathsf{P1}_2(\pt_1,q_1),\qquad
  \mathsf{P1}_3.
\]
\begin{enumerate}
\item In \(\mathsf{P1}_1\), on direct input
\[
  \pt_1\xrightarrow{\pt_2}(\alpha,1),
\]
it samples \(q_1\leftarrow\dexp\), sends to \(\Forw_1\)
\[
  (\alpha,3)\xrightarrow{((\alpha,4),(\pt_1,\pt_2,g\pow q_1))}(\alpha1,1),
\]
and moves to \(\mathsf{P1}_2(\pt_1,q_1)\).

\item In \(\mathsf{P1}_2(\pt_1,q_1)\), on a delivery from \(\Forw_2\) carrying
\(g\pow q_2\), it sends
\[
  (\alpha,1)\xrightarrow{g\pow(q_1q_2)}\pt_1
\]
and moves to \(\mathsf{P1}_3\).

\item In \(\mathsf{P1}_3\), all messages are rejected.
\end{enumerate}

Party~2 has two states:
\[
  \mathsf{P2}_1,\qquad
  \mathsf{P2}_2.
\]
\begin{enumerate}
\item In \(\mathsf{P2}_1\), on a delivery from \(\Forw_1\) carrying
\((\pt_1,\pt_2,g\pow q_1)\), it samples \(q_2\leftarrow\dexp\), sends
\[
  (\alpha,2)\xrightarrow{(\pt_1,g\pow(q_1q_2))}\pt_2,
\]
then sends to \(\Forw_2\)
\[
  (\alpha,4)\xrightarrow{((\alpha,3),g\pow q_2)}(\alpha2,1),
\]
and moves to \(\mathsf{P2}_2(q_2)\).

\item In \(\mathsf{P2}_2\), all messages are rejected.
\end{enumerate}

The internal distribution loop only routes legal messages among the parties and
the two forwarding sub-functionalities.  It performs no cryptographic
computation.
\end{definition}

\begin{definition}[The hybrid secure-message protocol \(\SMCReal(KE)\)]
For a key-exchange sub-functionality \(KE\), \(\SMCReal(KE)\) has address
\(\alpha\), adversary address \(\beta\), sub-functionality \(\Forw\) at
\(\alpha1\), sub-functionality \(KE\) at \(\alpha2\), and ports \(1,2,3,4\).
Ports \(1,3\) belong to Party~1; ports \(2,4\) belong to Party~2.

Party~1 has states
\[
  \mathsf{SP1}_1,\quad \mathsf{SP1}_2(\pt_1,\pt_2,t),\quad \mathsf{SP1}_3.
\]
\begin{enumerate}
\item In \(\mathsf{SP1}_1\), on
\[
  \pt_1\xrightarrow{(\pt_2,t)}(\alpha,1),
\]
it calls Party~1 of the KE sub-functionality:
\[
  (\alpha,3)\xrightarrow{(\alpha,4)}(\alpha2,1),
\]
stores \((\pt_1,\pt_2,t)\), and moves to
\(\mathsf{SP1}_2(\pt_1,\pt_2,t)\).

\item In \(\mathsf{SP1}_2(\pt_1,\pt_2,t)\), on
\[
  (\alpha2,1)\xrightarrow{k}(\alpha,3),
\]
it computes \(e=\inj(t)\mul k\), sends to \(\Forw\)
\[
  (\alpha,3)\xrightarrow{((\alpha,4),(\pt_1,\pt_2,e))}(\alpha1,1),
\]
and moves to \(\mathsf{SP1}_3\).

\item In \(\mathsf{SP1}_3\), all messages are rejected.
\end{enumerate}

Party~2 has states
\[
  \mathsf{SP2}_1,\quad \mathsf{SP2}_2(k),\quad \mathsf{SP2}_3.
\]
\begin{enumerate}
\item In \(\mathsf{SP2}_1\), on
\[
  (\alpha2,2)\xrightarrow{((\alpha,3),k)}(\alpha,4),
\]
it sends the second-phase KE activation
\[
  (\alpha,4)\xrightarrow{\star}(\alpha2,2),
\]
stores \(k\), and moves to \(\mathsf{SP2}_2(k)\).

\item In \(\mathsf{SP2}_2(k)\), on a delivery from \(\Forw\) carrying
\((\pt_1,\pt_2,e)\), it computes
\[
  t=\oget(\proj(e\mul\kinv(k))),
\]
sends
\[
  (\alpha,2)\xrightarrow{(\pt_1,t)}\pt_2,
\]
and moves to \(\mathsf{SP2}_3\).

\item In \(\mathsf{SP2}_3\), all messages are rejected.
\end{enumerate}
\end{definition}

\subsection{DDH games and the key-exchange reduction}

\begin{definition}[DDH games]
For an algorithm \(B:\Key^3\to\{0,1\}\), define
\[
\begin{array}{ll}
\DDHOne(B):&
q_1\leftarrow\dexp;\ q_2\leftarrow\dexp;\
\res\leftarrow B(g\pow q_1,g\pow q_2,g\pow(q_1q_2));\
\text{return }\res,\\[0.4em]
\DDHTwo(B):&
q_1\leftarrow\dexp;\ q_2\leftarrow\dexp;\ q_3\leftarrow\dexp;\
\res\leftarrow B(g\pow q_1,g\pow q_2,g\pow q_3);\
\text{return }\res.
\end{array}
\]
The DDH advantage of \(B\) at security parameter \(n\) is
\[
  \Delta_{\mathsf{DDH}}(B,n)
  =
  \left|
    \Pr[\DDHOne(B)(n)=1]
    -
    \Pr[\DDHTwo(B)(n)=1]
  \right|.
\]
\end{definition}

\begin{definition}[The KE simulator \(\KESim\)]
\(\KESim(A)\) wraps an adversary \(A\).  It forwards all environment/adversary
traffic not addressed to the simulated real functionality.  After the first
message from \(\KEIdeal\), it learns the ideal functionality address \(\alpha\)
and hence the real address it must spoof.

It internally simulates the two forwarding sub-functionalities of \(\KEReal\).
For the first forwarded DH message it samples \(q_1\leftarrow\dexp\) and exposes
\(g\pow q_1\) to \(A\) as \(\Forw_1\) would.  For the second forwarded DH
message it samples \(q_2\leftarrow\dexp\) and exposes \(g\pow q_2\) to \(A\) as
\(\Forw_2\) would.  Approvals obtained from \(A\) are translated into approvals
to \(\KEIdeal\).  The key output is not chosen by \(\KESim\); it is the
\(g\pow q\) sampled by \(\KEIdeal\).
\end{definition}

\begin{definition}[DDH adversary induced by \(Z,A\)]
For every environment \(Z\), adversary \(A\), addresses \(\alpha,\beta\), and
guard \(G\), define the DDH adversary
\[
  B_{Z,A,\alpha,\beta,G}(X,Y,W)
\]
as follows.  It runs the complete KE experiment with \(Z\) and \(A\), but the
three DH-related values are supplied externally:
\[
  g\pow q_1:=X,\qquad g\pow q_2:=Y,\qquad
  \text{agreed key}:=W.
\]
Concretely, the simulated first forwarder leaks \(X\), the simulated second
forwarder leaks \(Y\), Party~2's key output is \(W\), and Party~1's key output
is \(W\).  All other routing, adversarial scheduling, rejection behavior, and
state updates are exactly those of the EasyUC controller.  When the simulated
environment halts, \(B_{Z,A,\alpha,\beta,G}\) returns the environment's output
bit.
\end{definition}

\begin{lemma}[PPT closure for the DDH adversary]
If \(Z\) and \(A\) are PPT and the controller, message encoders, address tests,
group operations, and sampling from \(\dexp\) are PPT, then
\(B_{Z,A,\alpha,\beta,G}\) is PPT.
\end{lemma}

\begin{proof}
The machine \(B_{Z,A,\alpha,\beta,G}\) performs one execution of the same
controller used in the UC experiment.  It replaces the local instructions that
would sample \(q_1,q_2\) and compute \(g\pow q_1,g\pow q_2,g\pow(q_1q_2)\) by
reads of its three inputs \(X,Y,W\).  It does not add rewinding or unbounded
search.  The number of activations is bounded by the same polynomial bound as
the experiment with \(Z\) and \(A\).  Each additional operation is a fixed
number of tuple constructions and equality/address tests.  Therefore its
running time is bounded by the experiment's polynomial plus a fixed polynomial
overhead.
\end{proof}

\begin{lemma}[DDH game equality for real KE]
For all \(Z,A,\alpha,\beta,G,n\),
\[
\Pr[\Exec(\KEReal,A,Z;\alpha,\beta,G,n)=1]
=
\Pr[\DDHOne(B_{Z,A,\alpha,\beta,G})(n)=1].
\]
\end{lemma}

\begin{proof}
Couple the two executions by using the same coins for \(Z,A\), the controller,
and all non-DH randomized choices.  In \(\DDHOne\), the inputs to
\(B_{Z,A,\alpha,\beta,G}\) are
\[
  X=g\pow q_1,\qquad Y=g\pow q_2,\qquad W=g\pow(q_1q_2)
\]
for uniformly sampled \(q_1,q_2\).  These are exactly the values produced in
\(\KEReal\): \(\Forw_1\) leaks \(g\pow q_1\), \(\Forw_2\) leaks \(g\pow q_2\),
and both parties output \(g\pow(q_1q_2)\).  Every other transition is the same
controller transition with the same input and the same state.  Induction on the
activation number gives identical full traces, and therefore identical final
environment bits.  Taking probabilities gives the equality.
\end{proof}

\begin{lemma}[DDH game equality for ideal KE with simulator]
For all \(Z,A,\alpha,\beta,G,n\) such that \(2\notin G\),
\[
\Pr[\Exec(\KEIdeal,\KESim(A),Z;\alpha,\beta,G,n)=1]
=
\Pr[\DDHTwo(B_{Z,A,\alpha,\beta,G})(n)=1].
\]
\end{lemma}

\begin{proof}
In \(\DDHTwo\), the inputs to \(B_{Z,A,\alpha,\beta,G}\) are
\[
  X=g\pow q_1,\qquad Y=g\pow q_2,\qquad W=g\pow q_3,
\]
where \(q_1,q_2,q_3\) are mutually independent and uniform.  In
\(\Exec(\KEIdeal,\KESim(A),Z)\), the simulator independently samples
\(q_1,q_2\) and exposes \(g\pow q_1,g\pow q_2\) through the two simulated
forwarders, while \(\KEIdeal\) independently samples \(q_3\) and outputs
\(g\pow q_3\) as the session key.  The guard condition \(2\notin G\) ensures
that the environment cannot directly inject messages on the simulator port used
by \(\KESim\), so every message on that port arises only from \(\KEIdeal\) or
\(\KESim\).

Couple the two experiments by using the same \(q_1,q_2,q_3\), the same coins
for \(Z,A\), and the same controller coins.  The simulated forwarder messages
in \(\KESim(A)\) are identical to the forwarder messages generated by
\(B_{Z,A,\alpha,\beta,G}\) from \(X,Y\).  The key messages output by
\(\KEIdeal\) are identical to the key messages generated by
\(B_{Z,A,\alpha,\beta,G}\) from \(W\).  All approvals and non-KE messages are
routed identically.  Induction on activations yields identical traces and hence
identical final bits.
\end{proof}

\begin{theorem}[Concrete DDH security of key exchange]
For every PPT environment \(Z\), PPT adversary \(A\), valid addresses
\(\alpha,\beta\), input guard \(G\) with \(2\notin G\), and security parameter
\(n\),
\[
\begin{aligned}
&\left|
 \Pr[\Exec(\KEReal,A,Z;\alpha,\beta,G,n)=1]
 -
 \Pr[\Exec(\KEIdeal,\KESim(A),Z;\alpha,\beta,G,n)=1]
\right|\\
&\hspace{4em}
\leq
\Delta_{\mathsf{DDH}}(B_{Z,A,\alpha,\beta,G},n).
\end{aligned}
\]
\end{theorem}

\begin{proof}
By the real-game equality lemma,
\[
 \Pr[\Exec(\KEReal,A,Z;\alpha,\beta,G,n)=1]
 =
 \Pr[\DDHOne(B_{Z,A,\alpha,\beta,G})(n)=1].
\]
By the ideal-game equality lemma,
\[
 \Pr[\Exec(\KEIdeal,\KESim(A),Z;\alpha,\beta,G,n)=1]
 =
 \Pr[\DDHTwo(B_{Z,A,\alpha,\beta,G})(n)=1].
\]
Substituting these two equalities into the left-hand side gives exactly
\[
\left|
    \Pr[\DDHOne(B_{Z,A,\alpha,\beta,G})(n)=1]
    -
    \Pr[\DDHTwo(B_{Z,A,\alpha,\beta,G})(n)=1]
\right|
=
\Delta_{\mathsf{DDH}}(B_{Z,A,\alpha,\beta,G},n).
\]
\end{proof}

\subsection{SMC in the \(\KEIdeal\)-hybrid world}

\begin{definition}[The SMC simulator \(\SMCSim\)]
\(\SMCSim(A)\) wraps \(A\).  It forwards all traffic not addressed to the
simulated SMC functionality.  Upon receiving from \(\SMCIdeal\) the
adversarial notification
\[
  (\alpha,3)\xRightarrow{(\pt_1,\pt_2)}(\beta,3),
\]
it simulates the adversarial view of a run of \(\SMCReal(\KEIdeal)\):
\begin{enumerate}
\item it simulates the two adversarial notifications of the KE sub-functionality
at address \(\alpha2\), using port \(2\) of the adversary as in the KE
simulator convention;
\item it samples a fresh \(r\leftarrow\Key\) uniformly and uses \(r\) as the
simulated ciphertext sent through the forwarding sub-functionality at
\(\alpha1\);
\item it exposes to \(A\) the forwarding notification
\[
  (\alpha1,1)\xRightarrow{((\alpha,3),(\alpha,4),(\pt_1,\pt_2,r))}(\beta,1);
\]
\item when \(A\) approves the simulated forwarding request, it sends the
approval to \(\SMCIdeal\) on port \(3\).
\end{enumerate}
All messages sent by \(A\) to the simulated forwarding and KE sub-functionalities
are answered exactly as those sub-functionalities would answer, except that the
ciphertext is sampled as \(r\) and therefore does not depend on the plaintext.
\end{definition}

\begin{lemma}[Hybrid SMC ciphertext replacement]
For every fixed plaintext \(t\), every fixed environment/adversary history
\(h\) before the SMC ciphertext is generated, and every continuation \(C\),
\[
  \Pr_{k\leftarrow \Key}[C(h,\inj(t)\mul k)=1]
  =
  \Pr_{r\leftarrow \Key}[C(h,r)=1].
\]
\end{lemma}

\begin{proof}
By the one-time-pad masking lemma, \(\inj(t)\mul k\), with uniform \(k\), is
uniform over \(\Key\).  The random variable \(r\) is also uniform over
\(\Key\).  Since \(C(h,\cdot)\) is an arbitrary deterministic predicate after
conditioning on all coins outside the sampled key, equal input distributions
give equal output probabilities.  Averaging over the previously fixed history
preserves equality.
\end{proof}

\begin{lemma}[State correspondence for SMC hybrid and ideal]
Assume \(3\notin G\).  There is a relation \(R_{\mathsf{SMC}}\) between states
of the experiments
\[
  \Exec(\SMCReal(\KEIdeal),A,Z;\alpha,\beta,G,n)
\]
and
\[
  \Exec(\SMCIdeal,\SMCSim(A),Z;\alpha,\beta,G,n)
\]
with the following properties.
\begin{enumerate}
\item The initial states are related.
\item If related states have no pending SMC ciphertext, then a legal controller
step on either side is matched by zero or more legal controller steps on the
other side, preserving \(R_{\mathsf{SMC}}\) and producing the same visible
message.
\item If the hybrid side generates the ciphertext \(\inj(t)\mul k\), then the
ideal/simulator side generates \(r\leftarrow\Key\); replacing the former by the
latter preserves the distribution of every continuation.
\item If the hybrid receiver decrypts and sends \((\pt_1,t)\) to \(\pt_2\), the
ideal side first causes \(\SMCIdeal\) to send the same \((\pt_1,t)\) to
\(\pt_2\).  After this finite matching sequence the two states are again
related.
\end{enumerate}
\end{lemma}

\begin{proof}
Define \(R_{\mathsf{SMC}}\) by equality of the environment state, equality of
the wrapped adversary state, equality of the controller queue outside the
simulated SMC internals, equality of the stored external ports
\((\pt_1,\pt_2)\), equality of the KE control state up to the fact that in the
ideal experiment KE traffic is generated inside \(\SMCSim\), and equality of
the final message delivered to the external receiver whenever such a message is
pending.

The initial states are related because all local states are initialized to their
first state and all queues are empty.

For non-ciphertext steps, inspect the unique machine activated by the
controller.  If it is \(Z\) or the wrapped adversary \(A\), both experiments
call the same machine on the same visible message, so they return the same next
message and next local state.  If it is a routing step or a rejection step, the
message classification and guard predicates are the same on both sides.  If it
is a KE or forwarding notification, \(\SMCSim\) was defined to emit exactly the
notification that \(\SMCReal(\KEIdeal)\) would have emitted.  Thus each such
step is matched directly.

For the ciphertext step, the hybrid ciphertext is \(\inj(t)\mul k\) where the
key \(k\) is uniform from \(\KEIdeal\).  The ideal/simulator ciphertext is an
independent uniform \(r\).  The previous lemma shows that every continuation has
the same output distribution after this replacement.

For the final receiver step, the hybrid receiver computes
\[
  \oget(\proj((\inj(t)\mul k)\mul\kinv(k)))=t
\]
by one-time-pad correctness and outputs \((\pt_1,t)\) to \(\pt_2\).  The ideal
side, after the simulator's approval, causes \(\SMCIdeal\) to output the same
\((\pt_1,t)\) to \(\pt_2\).  The two-step ideal sequence is finite and restores
the relation.
\end{proof}

\begin{theorem}[Perfect SMC security in the \(\KEIdeal\)-hybrid]
For every PPT environment \(Z\), PPT adversary \(A\), valid addresses
\(\alpha,\beta\), guard \(G\) with \(3\notin G\), and security parameter \(n\),
\[
\Pr[\Exec(\SMCReal(\KEIdeal),A,Z;\alpha,\beta,G,n)=1]
=
\Pr[\Exec(\SMCIdeal,\SMCSim(A),Z;\alpha,\beta,G,n)=1].
\]
\end{theorem}

\begin{proof}
Use the state correspondence lemma and induct on the number of controller
activations until the environment halts.  The base case is initial-state
correspondence.  In the induction step, the correspondence lemma gives a
matching finite sequence on the other side and preserves the distribution of
the continuation.  The only randomized non-lockstep transition is ciphertext
generation, where equality of all continuations is exactly the ciphertext
replacement lemma.  Therefore the distribution of the final environment bit is
identical.
\end{proof}

\subsection{Bridging the SMC context to the KE experiment}

\begin{definition}[Composed environment]
For \(X\in\{\KEReal,\KEIdeal\}\), define
\[
  \CompEnv_X(Z)
\]
to be the environment that internally runs \(Z\) together with the outer SMC
context \(\SMCReal(\Box)\), where the hole \(\Box\) is implemented by a KE
stub.  The KE stub forwards calls to the external functionality \(X\) through
the surrounding interface.  An adversary stub forwards adversarial messages
between the internal SMC context and the external adversary.  If one stub
receives a message that must be returned through the other stub, it stores the
message in a shared one-message mailbox; the other stub returns the mailbox
contents at its next activation.  The lower guard \(G_{\ell}\) is the guard
seen by \(Z\); the upper guard \(G_h\) is the guard used by the external KE
experiment.
\end{definition}

\begin{lemma}[PPT closure for composed environments]
If \(Z\) is PPT, then \(\CompEnv_X(Z)\) is PPT for
\(X\in\{\KEReal,\KEIdeal\}\).
\end{lemma}

\begin{proof}
\(\CompEnv_X(Z)\) runs one copy of \(Z\), one copy of the finite-state SMC
context, and two finite-state stubs.  Each activation performs one activation
of one of these machines plus a constant number of mailbox and routing
operations.  The SMC context contains only the one-shot \(\Forw\) and one call
to \(X\); its distribution loops are equipped with the same decreasing
termination metrics as the original EasyUC machines.  Thus the number of
internal routing steps per external activation is bounded by a fixed
polynomial.  Therefore a polynomial bound for \(Z\) gives a polynomial bound
for \(\CompEnv_X(Z)\).
\end{proof}

\begin{lemma}[Bridge for \(\KEReal\)]
Let \(G_{\ell}\subseteq G_h\) and \(1\in G_h\).  For every \(Z,A,\alpha,\beta,n\),
\[
\Pr[\Exec(\SMCReal(\KEReal),A,Z;\alpha,\beta,G_{\ell},n)=1]
=
\Pr[\Exec(\KEReal,A,\CompEnv_{\KEReal}(Z);\alpha2,\beta,G_h,n)=1].
\]
\end{lemma}

\begin{proof}
Construct a trace map from the left experiment to the right experiment.  The
right experiment runs the same \(Z\), the same \(A\), and the same SMC context,
but the KE sub-functionality is exposed as the external functionality at
address \(\alpha2\).  All calls from the SMC context to KE are sent through the
KE stub; all returned messages are put back into the SMC context through the
stub.  All adversarial messages that the SMC context sends to \(A\) are routed
through the adversary stub.  The assumption \(G_{\ell}\subseteq G_h\) ensures
that every adversarial message allowed in the left experiment is allowed in the
right experiment, and \(1\in G_h\) permits forwarding-control traffic.

The map preserves the local state of \(Z\), \(A\), \(\SMCReal\)'s two parties,
\(\Forw\), \(\KEReal\)'s two parties, and \(\KEReal\)'s two forwarding
sub-functionalities.  It also records whether a message is currently in the
stub mailbox.  Induct on the lexicographic pair
\[
  (N,M),
\]
where \(N\) is the remaining external activation budget and \(M\) is the sum of
the termination metrics of the currently active distribution loops.  A direct
environment or adversary step decreases \(N\).  An internal routing step either
decreases \(M\) or returns no message and leaves the visible state unchanged.
The transition definitions of the stubs make the mailbox cases inverse to the
corresponding direct routing cases.  Hence every left trace has a unique right
trace with the same probability and the same final environment bit.  The same
construction in reverse gives uniqueness in the other direction.  Therefore the
two Bernoulli distributions are equal.
\end{proof}

\begin{lemma}[Bridge for \(\KEIdeal\)]
Let \(G_{\ell}\subseteq G_h\) and \(1\in G_h\).  For every \(Z,A,\alpha,\beta,n\),
\[
\Pr[\Exec(\SMCReal(\KEIdeal),A,Z;\alpha,\beta,G_{\ell},n)=1]
=
\Pr[\Exec(\KEIdeal,A,\CompEnv_{\KEIdeal}(Z);\alpha2,\beta,G_h,n)=1].
\]
\end{lemma}

\begin{proof}
The proof is identical to the proof of the bridge for \(\KEReal\), replacing
\(\KEReal\)'s local state and two forwarding sub-functionalities by
\(\KEIdeal\)'s five-state machine.  The termination metric is the remaining
number of states before \(\mathsf{KState}_5\), plus the finite routing-loop
metric.  Each accepted input either advances the five-state machine or returns
no output; no transition increases the metric.  The same trace bijection and
probability preservation argument applies.
\end{proof}

\subsection{Full SMC security}

\begin{definition}[Composed SMC simulator]
For an adversary \(A\), define
\[
  \SMCSimComp(A)=\SMCSim(\KESim(A)).
\]
The simulator ports are distinct: port \(2\) is reserved for \(\KESim\)'s
communication with \(\KEIdeal\), and port \(3\) is reserved for
\(\SMCSim\)'s communication with \(\SMCIdeal\).
\end{definition}

\begin{theorem}[Concrete UC security of real SMC]
Let \(G_{\ell}\) be an input guard such that \(2\notin G_{\ell}\) and
\(3\notin G_{\ell}\).  Let
\[
  G_h = G_{\ell}\cup\{1\}.
\]
For every PPT environment \(Z\), PPT adversary \(A\), valid addresses
\(\alpha,\beta\), and security parameter \(n\), let
\[
  B =
  B_{\CompEnv_{\KEReal}(Z),A,\alpha2,\beta,G_h}.
\]
Then
\[
\begin{aligned}
&\left|
 \Pr[\Exec(\SMCReal(\KEReal),A,Z;\alpha,\beta,G_{\ell},n)=1]
 -
 \Pr[\Exec(\SMCIdeal,\SMCSimComp(A),Z;\alpha,\beta,G_{\ell},n)=1]
\right|\\
&\hspace{4em}
\leq \Delta_{\mathsf{DDH}}(B,n).
\end{aligned}
\]
\end{theorem}

\begin{proof}
First apply the bridge for \(\KEReal\):
\[
\Pr[\Exec(\SMCReal(\KEReal),A,Z;\alpha,\beta,G_{\ell},n)=1]
=
\Pr[\Exec(\KEReal,A,\CompEnv_{\KEReal}(Z);\alpha2,\beta,G_h,n)=1].
\]
Since \(2\notin G_{\ell}\) and \(G_h=G_{\ell}\cup\{1\}\), we have
\(2\notin G_h\).  Hence the concrete DDH security theorem for key exchange
applies to the right-hand side and gives
\[
\begin{aligned}
&\left|
\Pr[\Exec(\KEReal,A,\CompEnv_{\KEReal}(Z);\alpha2,\beta,G_h,n)=1]
-
\Pr[\Exec(\KEIdeal,\KESim(A),\CompEnv_{\KEReal}(Z);\alpha2,\beta,G_h,n)=1]
\right|\\
&\hspace{4em}\leq \Delta_{\mathsf{DDH}}(B,n).
\end{aligned}
\]
The composed environment used with \(\KEIdeal\) is extensionally the same
machine as \(\CompEnv_{\KEIdeal}(Z)\) after replacing the external
functionality by \(\KEIdeal\); therefore
\[
\Pr[\Exec(\KEIdeal,\KESim(A),\CompEnv_{\KEReal}(Z);\alpha2,\beta,G_h,n)=1]
=
\Pr[\Exec(\KEIdeal,\KESim(A),\CompEnv_{\KEIdeal}(Z);\alpha2,\beta,G_h,n)=1].
\]
Now apply the bridge for \(\KEIdeal\), with adversary \(\KESim(A)\), in the
reverse direction:
\[
\Pr[\Exec(\KEIdeal,\KESim(A),\CompEnv_{\KEIdeal}(Z);\alpha2,\beta,G_h,n)=1]
=
\Pr[\Exec(\SMCReal(\KEIdeal),\KESim(A),Z;\alpha,\beta,G_{\ell},n)=1].
\]
Finally apply perfect SMC security in the \(\KEIdeal\)-hybrid, with adversary
\(\KESim(A)\).  Since \(3\notin G_{\ell}\),
\[
\Pr[\Exec(\SMCReal(\KEIdeal),\KESim(A),Z;\alpha,\beta,G_{\ell},n)=1]
=
\Pr[\Exec(\SMCIdeal,\SMCSim(\KESim(A)),Z;\alpha,\beta,G_{\ell},n)=1].
\]
Combining the displayed equalities with the DDH inequality yields exactly the
claimed bound, and \(\SMCSim(\KESim(A))=\SMCSimComp(A)\) by definition.
\end{proof}

\begin{corollary}[Asymptotic UC security of real SMC]
If the DDH assumption holds for the group family, then
\[
  \SMCReal(\KEReal)
\]
UC-realizes \(\SMCIdeal\) with simulator \(\SMCSimComp\), for guards excluding
the simulator ports \(2\) and \(3\).
\end{corollary}

\begin{proof}
By the PPT closure lemmas, \(B\) in the concrete SMC theorem is a PPT DDH
adversary whenever \(Z\) and \(A\) are PPT.  The DDH assumption says that
\(\Delta_{\mathsf{DDH}}(B,n)\) is negligible for every such \(B\).  The concrete
SMC theorem upper-bounds the distinguishing advantage of \(Z\) by this
negligible function.  Therefore the distinguishing advantage is negligible for
all PPT \(Z,A\), which is precisely UC realization with simulator
\(\SMCSimComp\).
\end{proof}